\clearpage
\begingroup
\NoSourceHead
\HistoricalFoot{XI1.}{31}
\SourceLeaf{246}
\begin{center}
[III.]\par\medskip
{\large [ZUR BAHNBESTIMMUNG DER CERES.]}\par\medskip
\rule{.44\linewidth}{.3pt}\\
{\small [Aus Scheda Ah, S. 1—8, 10—13.]}\\
\rule{.44\linewidth}{.3pt}
\end{center}
\noindent{\small [S. 1]}\par
\begingroup\raggedright\noindent\hspace*{\parindent}\textls[100]{Berechnung der Bahn des neuen von \textsc{Piazzi} entdeckten Planeten aus folgenden drei, von den Wirkungen der Aberration, Präcession der Nachtgleichen und den wahrscheinlichsten Beobachtungsfehlern bereits befreieten Beobachtungen.}\par\endgroup
\begin{center}
\textls[100]{[Ableitung der III. Elemente nach der zweiten Verbesserungsmethode.]}
\end{center}
{\footnotesize\setlength{\arraycolsep}{2.5pt}
\[
\begin{array}{rlllll}
&\left[\begin{gathered}\text{Mittlere}\\\text{Sonnenzeit}\end{gathered}\right]
&\left[\begin{gathered}\text{Geocentri-}\\\text{sche Länge}\end{gathered}\right]
&\left[\begin{gathered}\text{Geocentrische}\\\text{Breite, südlich}\end{gathered}\right]
&\left[\begin{gathered}\text{Länge d. Sonne}\\+20''\text{ Aberration}\end{gathered}\right]
&\left[\begin{gathered}\text{Log. Abstand der}\\\text{Sonne von d. Erde}\end{gathered}\right]\\[3pt]
\text{Januar}&1{,}363\,4005&53^0\,22'\,59''9&3^0\;6'\,40''8&281^0\;1'\,30''90&9{,}992\,6156\\
&21{,}308\,3647&53\;34\;21''9&1\;46\;7''0&301\;20\;37''55&9{,}993\,1434\\
&42{,}258\,3125&56\;26\;37''9&0\;35\;56''5&322\;35\;5''76&9{,}994\,5823.
\end{array}
\]}
\noindent\hspace*{\parindent}\textbf{Erste Hypothese.}\quad Ort des $\AstNode=81^0\,2'\,35''$; $i=10^0\,36'\,30''$\textsuperscript{[*)]}.\begingroup\renewcommand{\thefootnote}{[*)}\footnotetext{Diese Werte entsprechen den II. Elementen; vgl. Werke VI, S. 200.]}\endgroup
\par\medskip
\noindent\hspace*{\parindent}\textbf{Berechnung der ersten Beobachtung.}
\par\medskip\noindent [Formeln:
\[
\frac R\Delta=\frac{\sin(\alpha-\AstNode)-\frac\theta g}{\sin(V-\AstNode)},
\]
\[
\frac r\Delta=\surd\left(1+\theta\theta+\frac{RR}{\Delta\Delta}-2\frac R\Delta\cos(V-\alpha)\right),
\]
\[
\operatorname{tang}(v-\AstNode)=\frac{\frac\theta{\sin i}\operatorname{tang}(V-\AstNode)}
{\frac\theta g-\frac{\sin(\alpha-V)}{\cos(V-\AstNode)}},
\]
\begin{center}\textbf{Zur Probe:}\end{center}
\[
\sin(v-\AstNode)=\frac\Delta r\mathbin{.}\frac\theta{\sin i}],
\]
\clearpage
\HistoricalHead{242}{L}{THEORETISCHE ASTRONOMIE. NACHLASS.}
\SourceLeaf{247}
\[
\begin{array}{rl}
V-\alpha=&227^0\,38'\,31''00\\
V-\AstNode=&199\;58\;55''90\\
\alpha-\AstNode=&-27\;39\;35''10,
\end{array}
\]
\[
\begin{array}{rl}
\log\theta=&8{,}735\,2511\,[n]\\
\log g=&9{,}272\,5272\\\cline{2-2}
{}\left[\log\dfrac\theta g=\right]&9{,}462\,7239\,[n]
\end{array}
\]
\[
\begin{array}{rl}
\sin(\alpha-\AstNode)=[-]&0{,}464\,21995\\
{}\left[\dfrac\theta g=-\right]&0{,}290\,21772\\\cline{2-2}
{}[\text{Differenz}=-]&0{,}174\,00223
\end{array}
\]
\[
\begin{array}{rl}
{}\left[\log\left(\sin(\alpha-\AstNode)-\dfrac\theta g\right)=\right]&9{,}240\,5548\,[n]\\[3pt]
{}[\log\sin(V-\AstNode)=]&9{,}533\,6807\,[n]\\\cline{2-2}
{}\left[\log\dfrac R\Delta=\right]&9{,}706\,8741\\[3pt]
{}[\log R=]&9{,}992\,6156\\\cline{2-2}
{}[\log\Delta=]&0{,}285\,7415
\end{array}
\]
\noindent\rule{\linewidth}{.3pt}
\[
\begin{array}{rl}
{}\left[\log2\dfrac R\Delta=\right]&0{,}007\,9041\\[3pt]
{}[\log\cos(\alpha-V)=]&9{,}828\,5063\,[n]\\\cline{2-2}
{}[\text{Summe}=]&9{,}836\,4104\,[n]
\end{array}
\]
\[
\begin{array}{rl}
{}[1+\theta\theta=]&1{,}002\,95462\\[8pt]
{}\left[\dfrac{RR}{\Delta\Delta}=\right]&0{,}259\,26756\\[8pt]
{}\left[-2\dfrac R\Delta\cos(\alpha-V)=\right]&0{,}686\,13635\\[8pt]\cline{2-2}
{}\left[\dfrac{rr}{\Delta\Delta}=\right]&1{,}948\,35853\\[8pt]
{}\left[\log\dfrac{rr}{\Delta\Delta}=\right]&0{,}289\,6688\\[8pt]
{}\left[\log\dfrac r\Delta=\right]&0{,}144\,8344\\[8pt]
{}[\log\Delta=]&0{,}285\,7415\\[8pt]\cline{2-2}
\log r=&0{,}430\,5759
\end{array}
\]
\noindent\rule{\linewidth}{.3pt}
\[
\begin{array}{rl}
\sin(\alpha-V)\;\text{. . . .}&9{,}868\,6143\\
\cos(V-\AstNode)\;\text{. . . .}&9{,}973\,0349\,[n]\\\cline{2-2}
{}\left[-\dfrac{\sin(\alpha-V)}{\cos(V-\AstNode)}\;\text{. . . .}\right]&9{,}895\,5794\\[3pt]
{}[\text{Numerus}=]&0{,}786\,28393\\[3pt]
{}\left[\dfrac\theta g=-\right]&0{,}290\,21772\\\cline{2-2}
{}\left[\dfrac\theta g-\dfrac{\sin(\alpha-V)}{\cos(V-\AstNode)}=\right]&0{,}496\,06621
\end{array}
\]
\[
\begin{array}{rl}
\operatorname{tang}(V-\AstNode)\;\text{. . . .}&9{,}560\,6458\\[3pt]
{}\left[\mathrm{C.}\log\left(\dfrac\theta g-\dfrac{\sin(\alpha-V)}{\cos(V-\AstNode)}\right)=\right]&0{,}304\,4603\\[3pt]
{}\left[\dfrac\theta{\sin i}\;\text{. . . .}\right]&9{,}470\,2108\,[n]\\\cline{2-2}
{}[\operatorname{tg}(v-\AstNode)\;\text{. . . .}]&9{,}335\,3169\,[n]\\[3pt]
{}[\cos(v-\AstNode)\;\text{. . . .}]&9{,}990\,0595\\\cline{2-2}
{}[\sin(v-\AstNode)\;\text{. . . .}]&9{,}325\,3764\,[n]\\[5pt]
{}[v-\AstNode=-]&12^0\,12'\,43''76\\
{}[\AstNode=]&81\;2\;35''00\\\cline{2-2}
v=&68\;49\;51''24.
\end{array}
\]
\clearpage
\HistoricalHead{243}{R}{ZUR BAHNBESTIMMUNG DER CERES.}
\SourceLeaf{248}
\noindent\hspace*{\parindent}\textbf{[Zur Probe:]}
\[
\begin{array}{rl}
{}\left[\dfrac\theta{\sin i}\;\text{. . . .}\right]&9{,}470\,2108\,[n]\\[3pt]
{}\left[\dfrac r\Delta\;\text{. . . .}\right]&0{,}144\,8344\\\cline{2-2}
{}[\sin(v-\AstNode)\;\text{. . . .}]&9{,}325\,3764\,[n].
\end{array}
\]
\noindent\hspace*{\parindent}\textbf{Zweite Beobachtung.}
\[
\begin{array}{rl}
V-\alpha=&247^0\,46'\,15''65\\
V-\AstNode=&220\;18\;2''55\\
\alpha-\AstNode=&-27\;28\;13''10
\end{array}
\quad\text{[Die Rechnung ergab:]}\quad
\begin{array}{rl}
{}[\log\Delta=]&0{,}331\,9954\\
{}[v=]&73^0\,19'\,11''47\\
\log r=&0{,}428\,2441.
\end{array}
\]
\noindent{\small [S. 2]}\par
\noindent\hspace*{\parindent}\textbf{Dritte Beobachtung.}
\[
\begin{array}{rl}
V-\alpha=&266^0\;8'\,27''86\\
V-\AstNode=&241\;32\;30''76\\
\AstNode-\alpha=&24\;35\;57''10
\end{array}
\quad\text{[Die Rechnung ergab:]}\quad
\begin{array}{rl}
{}[\log\Delta=]&0{,}381\,8140\\
{}[v=]&78^0\,5'\,59''16\\
\log r=&0{,}425\,5949.
\end{array}
\]
\begin{center}\rule{.20\linewidth}{.3pt}\end{center}
\clearpage
\endgroup
